\documentclass[../../main.tex]{subfiles}

\begin{document}

\subsection{Calculations}

\begin{lemma}\label{lem:rational_sextic_ten_nodes}
    Let $C \subset \PP^2$ be an irreducible sextic curve with at worst $A_1$ singularities. Then $C$ is a rational curve if and only if it has exactly ten $A_1$ singularities.
\end{lemma}

\begin{proof}
    One has the genus-degree formula for singular plane curves:
    \[
    p_a(C) = {1\over 2}(d-1)(d-2) - \sum_p \delta_p
    ,\]
    where the sum is over the singular locus of $C$. Since $C$ has at worst $A_1$ singularities, each singularity contributes $\delta_p = 1$. Noting that $d=6$, we obtain the geometric genus $g = 10 - k$, where $k$ is the number of $A_1$ singularities. Thus $g = 0$ if and only if $k = 10$.
\end{proof}

\begin{remark}[Severi varieties and degenerate sextics]\label{rmk:severi-sextics}
    We note that an arbitrary rational irreducible sextic need not have exactly ten singularities, nor must its singularities be $A_1$ singularities. The condition $g=0$ merely requires that the sum of the $\delta$-invariants of its singularities equals 10. The locus of rational $A_1$-singular sextics forms a Severi variety $V_{6,10}$ \cite{BHO+11} whose generic element has precisely 10 $A_1$ singularities. However, $V_{6,10}$ contains lower-dimensional strata parametrizing rational sextics with fewer singularities of higher multiplicity. For instance, there exist rational $A_2$-singular sextics lacking 10 $A_1$ singularities entirely \cite{GK20}, provided their singularities satisfy the bounded condition $\sum \delta_p = 10$. The classical unnodal Coble surface construction specifically requires a sextic from the generic, 10-$A_1$ locus.
\end{remark}

\begin{remark}[Canonical class of Cobles]
Let $W\subset \PP^2$ be a generic irreducible rational sextic with ten $A_1$ singularities. The classical Coble surfaces $S$, first constructed in \cite{Cob19} in a study of Cremona transformations of $\PP^2$ which preserve such a rational sextic $W$. These surfaces can be realized as the blowup of $\PP^2$ along the divisor $D = p_1 + \cdots + p_{10}$ comprised of the ten $A_1$ singularities of $W$. In this construction, letting $C$ be the proper transform of $W$, one can show that $C$ is a rational curve and $\abs{-2K_S} = \ts{C}$ by adjunction. Moreover, the square of the canonical class can be computed as follows:
\begin{align*}
    K_S^2 &= \qty{ \pi^* K_{\PP^2} + \sum_{k=1}^{10}E_k}^2 \\
    &=\qty{\pi^* K_{\PP^2}}^2 + \qty{\sum_{k=1}^{10}E_k}^2 + 2\qty{ \sum_{k=1}^{10}\pi^* K_{\PP^2}\cdot E_k } \\
    &=\qty{-3H }^2 + \sum_{k=1}^{10}E_k^2  \\
    &= 9 + 10\cdot(-1) \\
    &= -1.
\end{align*}


Thus $S$ is of K3 type, and the 2-to-1 cover $X\to S$ branched along $C$ is a K3 surface with $\rho(Y) = 11$, where a primitive sublattice $\mathrm{I}_{1, 10}(2)$ of $\Pic(X)$ is generated by the pullbacks of a hyperplane class $E_0$ and ten exceptional curves $E_i$. We note that the group of deck transformations of $X\to S$ is generated by a nonsymplectic involution $\iota$ which preserves $\NS(X)$. 
Thus if $S$ is of K3 type, then the universal double cover $f: X\to S$ is a singular normal K3 surface $X$ at worst $A_1$ singularities which is naturally equipped with a nonsymplectic involution $\iota$. Moreover, by the proof of \cite[Lem. 2.2]{AD19}, a Hurwitz-type formula yields
\[
0 = 2K_X = 2\qty{ f^*(K_S) + \sum_{k=1}^{10} E_k } = f^*(2K_S + C)
,\]
yielding $C\sim -2K_S$.
\end{remark}

\begin{remark}[Invariants of antibicanonical curves]
By \cite[\S 3.1]{CD12}, if $S$ is Coble surface of K3 type with $n=1$ boundary components and $C \in \abs{-2K_S}$ is an irreducible curve, then
\[
p_{a}(C)=1+\frac{1}{2}\left(C^{2}+C \cdot K_{S}\right)=1+K_{S}^{2}=0,
\]
noting that $K_S^2 = -1$.
This forces $C$ to be a smooth rational curve with $C^2 = 4K_S^2 = -4$.
\end{remark}

\subsection{Coxeter Diagrams}

Conventions: 

\begin{enumerate}
    \item White vertices: $v^2 = -2$.
    \item Black vertices: $v^2 = -4$.
\end{enumerate}

\subsection{Classical and affine Dynkin diagrams}
\Cref{tab:dynkin-diagrams-table} is a table of (labeled) classical and affine Dynkin diagrams.

\begin{table}[H]
\centering
\resizebox{\textwidth}{!}{%
\begin{tabular}{@{}llll@{}}
\toprule
\multicolumn{2}{l}{Classical Type}          & \multicolumn{2}{l}{Affine Type}                              \\ \midrule
$A_n$ & \dynkin[label, labels={1,2,n-1, n}, edge length=.75cm]A{}   & $\tilde A_n$ & \dynkin[extended, label, labels={0,1,2,n-1, n}, edge length=.75cm]A{}        \\
$B_n$ & \dynkin[label, labels={1,2,,n-1, n}, edge length=.75cm]B{}  & $\tilde B_n$ & \dynkin[extended, label, labels={0,1,2,3,n-2, n-1, n}, edge length=.75cm]B{} \\
$C_n$ & \dynkin[label, labels={1,2,, n-1, n}, edge length=.75cm]C{} & $\tilde C_n$ & \dynkin[extended, label, labels={0,1,2,n-2, n-1, n}, edge length=.75cm]C{}   \\
$D_n$ & \dynkin[label, labels={1,2,,,n-1,n}, edge length=.75cm]D{}  & $\tilde D_n$ & \dynkin[extended, label, labels={0,1,2,3,,,n-1,n}, edge length=.75cm]D{}     \\
$E_6$ & \dynkin[label, edge length=.75cm]E6 & $\tilde E_6$ & \dynkin[extended, label, edge length=.75cm]E6 \\
$E_7$ & \dynkin[label, edge length=.75cm]E7 & $\tilde E_7$ & \dynkin[extended, label, edge length=.75cm]E7 \\
$E_8$ & \dynkin[label, edge length=.75cm]E8 & $\tilde E_8$ & \dynkin[extended, label, edge length=.75cm]E8 \\
$F_4$ & \dynkin[label, edge length=.75cm]F4 & $\tilde F_4$ & \dynkin[extended, label, edge length=.75cm]F4 \\
$G_2$ & \dynkin[label, edge length=.75cm]G2 & $\tilde G_2$ & \dynkin[extended, label, edge length=.75cm]G2 \\ \bottomrule
\end{tabular}%
}
\caption{Classical and Affine Dynkin Diagrams}
\label{tab:dynkin-diagrams-table}
\end{table}

\subsection{Formulae}

\begin{remark}
    Let $f: Y = \Bl_V X \birational X$ be the blowup of $X$ along a subvariety $V$. Then one has the adjunction formula
    \[
    K_Y = f^* K_X + (\codim_V(X) - 1)E
    .\]
    Note that if $X$ is a surface and $V$ is a divisor, this recovers 
    \[
    K_Y = f^* K_X + E
    .\]
\end{remark}

\begin{remark}
    Let $f: Y\to X$ be a branched cover with branch locus $B$ and ramification locus $R$. Then
    \[
    K_Y = f^* K_X + R, \qquad R \da \sum_{p\in X} \mathrm{len}\qty{ \Omega_{X/Y, p}} [p]
    .\]
\end{remark}

\subsection{Cusps}

\begin{table}[H]
\centering
\begin{tabular}{@{}llll@{}}
\toprule
Sterk Cusp & Vector                           & $\mathrm{div}_{T_\En}$ & $\mathrm{div}_{T_{\mathrm{K3} } }$ \\ \midrule
1          & $e$                              & 1                      & 1                                  \\
2          & $e'$                             & 2                      & 2                                  \\
3          & $e' + f' + \overline{\alpha}_8$  & 2                      & 1                                  \\
4          & $2e' + f' + \overline{\alpha}_1$ & 2                      & 1                                  \\
5          & $2e + 2f + \overline{\alpha}_1$  & 2                      & 1                                  \\ \bottomrule
\end{tabular}
\caption{Isotropic vectors in $F_{\En, 2}$ and their divisibilities.}
\label{tab:sterk-cusps}
\end{table}

\subsection{Coble lattices}

The following can be found in an unpublished note at \href{https://homepage.mi-ras.ru/~prokhoro/conf/isk20/Dolgachev.pdf}{https://homepage.mi-ras.ru/~prokhoro/conf/isk20/Dolgachev.pdf}.
Write $\abs{-2K_S} = \ts{C}$ where $C = C_1 + \cdots + C_n$ has $n$ irreducible components. By adjunction, $C_i\cong \PP^1$ and $C_i^2 = -4$, so $K_S^2 = -n$. Let $\Sigma$ be a general Coble set of points in $\PP^2$. There are 10 irreducible families of K3s obtained as the double cover of $S$ branched over $C$, where the fixed point locus of the deck transformation consists of $n$ smooth rational curves. Generally the rank of $\Pic(X)$ is $r=10+n$ and $\ell = 12-n$, and $K_S^2 = 9-\abs\Sigma = -n$.

\begin{tabular}{|c|c|c|c|c|c|}
\hline$n$ & $\abs\Sigma$ & $K_{\mathrm{V}}^{2}$ & $M = (r, l, \delta)$ & 2-elementary lattice $M$ & $N=M^{\perp}$ \\
\hline 1 & 10 & -1 & $(11,11,1)$ & $\mathrm{E}_{10}(2) \oplus \mathrm{A}_{1}$ & $\mathrm{I}^{2,9}(2)$ \\
2 & 11 & -2 & $(12,10,1)$ & $\mathrm{E}_{8}(2) \oplus \mathrm{U} \oplus \mathrm{A}_1^{\oplus 2}$ & $\mathrm{I}^{2,8}(2)$ \\
3 & 12 & -3 & $(13,9,1)$ & $\mathrm{D}_{4}^{\oplus 2} \oplus \mathrm{A}_{1}^{\oplus 3} \oplus \mathrm{U}(2)$ & $\mathrm{I}^{2,7}(2)$ \\
4 & 13 & -4 & $(14,8,1)$ & $\mathrm{D}_{4}^{\oplus 2} \oplus \mathrm{A}_{1}^{\oplus 4} \oplus \mathrm{U}(2)$ & $\mathrm{I}^{2,6}(2)$ \\
5 & 14 & -5 & $(15,7,1)$ & $\mathrm{E}_{8} \oplus \mathrm{A}_{1}^{\oplus 5} \oplus \mathrm{U}$ & $\mathrm{I}^{2,5}(2)$ \\
6 & 15 & -6 & $(16,6,1)$ & $\mathrm{E}_{10} \oplus \mathrm{A}_{1}^{\oplus 6}$ & $\mathrm{I}^{2,4}(2)$ \\
7 & 16 & -7 & $(17,5,1)$ & $\mathrm{E}_{8} \oplus \mathrm{D}_{6} \oplus \mathrm{A}_{1} \oplus \mathrm{U}(2)$ & $\mathrm{I}^{2,3}(2)$ \\
8 & 17 & -8 & $(18,4,0)$ & $\mathrm{E}_{8} \oplus \mathrm{D}_{8} \oplus \mathrm{U}(2)$ & $\mathrm{U}(2)^{\oplus 2}$ \\
8 & 17 & -8 & $(18,4,1)$ & $\mathrm{E}_{10} \oplus \mathrm{D}_{6} \oplus \mathrm{A}_{1}^{\oplus 2}$ & $\mathrm{I}^{2,2}(2)$ \\
9 & 18 & -9 & $(19,3,1)$ & $\mathrm{E}_{10} \oplus \mathrm{D}_{8} \oplus \mathrm{A}_{1}$ & $\mathrm{I}^{2,1}(2)$ \\
10 & 19 & -10 & $(20,2,1)$ & $\mathrm{E}_{10} \oplus \mathrm{D}_{10}$ & $(2\rangle^{\oplus 2}$ \\
\hline
\end{tabular}

This table is reproduced in \cite[Table 5.1, p.553]{CD12}. The fixed point locus is described by \cite[Eqn. 5.3.1]{CD12}:

$$
X^{g}=\left\{\begin{array}{ll}
\emptyset & \text { if }(r, l, \delta)=(10,10,0), \\
C_{1}^{(1)}+C_{2}^{(1)} & \text { if }(r, l, \delta)=(10,8,0), \\
C^{(g)}+\sum_{i=1}^{k} R_{i} & \text { otherwise }
\end{array}\right.
$$
where $C^{(g)}$ denotes a curve of genus $g \geq 0, R_{i}$ are disjoint (-2)-curves, and
$$
g=\frac{1}{2}(22-r-l), \quad k=\frac{1}{2}(r-l) .
$$
The ramification divisor for the involution on the K3 cover is described by $g=0$ and $k=n-1$ by \cite[Def. 5.4.3]{CD12}.

\subsection{Diagrams}
\hfill\newline
\begin{figure}[H]
    \centering
    \input{figures/tikzcd/fig_999_Appendix_1.tex}
    \caption{Mirror move algorithm 1}
    \label{fig:mirror-move-1}
\end{figure}

\begin{figure}[H]
    \centering
    \input{figures/tikzcd/fig_999_Appendix_2.tex}
    \caption{Mirror move algorithm 2}
    \label{fig:mirror-move-2}
\end{figure}

\hrule

% Coxeter diagram 7,7,1 and 9,9,1
\begin{figure}[H]
\centering
\input{figures/tikz/fig_999_Appendix_1.tex}
    \caption{The Coxeter diagrams $G_{(9,9,1)_1} = G_{\gens{2} \oplus E_8(2)}$ and $G_{(7,7,1)_1} = G_{A_1^{\oplus 7}}$ respectively.}
    \label{fig:coble-coxeter-diagrams}
\end{figure}

\subsection{Nodal Enriques Surfaces}

The lattices for $S_{\Nod}$ comes from \cite[Prop. 3.1]{DK13}. There is a primitive embedding
\begin{align*}
    S_{\En} &\injects S_{\Nod} \\
    U(2) \oplus E_{8}(2) &\injects U \oplus \gens{-4}  \oplus E_8(2) \\
    ( (e, f), x) &\mapsto ( (2\tilde f + \tilde g + h, \tilde g), x)
\end{align*}
where $U(2) = \gens{e,f}$, $U = \gens{\tilde e, \tilde f}$, and $\gens{-4} = \gens{h}$; this map is the identity on the $E_8(2)$ component. This formula is given on \cite[p.\, 561]{CD12}. There is is claimed that this induces a map $F_{S_{\Nod}} \injects F_{S_{\En}}$ where the former coincides with the moduli space of nodal Enriques surfaces.

\begin{align*}
S_{\Nod} &= (11, 10, 1)_1 = \gens{-4} \oplus E_{10}(2) & T_{\Nod} &= (11, 8, 1)_2 \,\,\,= \gens{4} \oplus U \oplus E_8(2)
\end{align*}

\end{document}